% disclaimers.tex — the RWI disclaimer catalogue, rendered (two-column build).
% Identical content to the one-column template; kept as a separate copy so
% each template dir stays self-contained (see design doc Part 3, "why two
% copies not one shared include"). IDs are the semantic ids of the single
% disclaimer catalogue (Chair Ruling v1 §A2/§C2) -- no numeric D1-D12 scheme.


\begin{disclaimers}
  \item[D-STANDPOINT] See \Cref{sec:standpoint}. [FILL: restate in one
    sentence for a reader who skips to Limitations.]
  \item[D-NONEXPERT] The author writes from the standpoint declared in
    \Cref{sec:standpoint} and is \emph{not} a substitute for a credentialed
    expert in any discipline this paper touches. [FILL: name the specific
    disciplines this paper draws on, and the disciplines it explicitly does
    not claim competence in --- see the Standpoint paragraph's "disciplines
    not claimed" list.]
  \item[D-SCOPE] [FILL if EMPIRICAL genre: n, population, sampling method,
    and the boundary past which findings should not be generalized. Else:
    not applicable to this genre.]
  \item[D-NONCLAIM] [FILL: an explicit "we do NOT claim ..." ledger. List at
    least the claims a reader might mistakenly infer from the title or
    abstract that this paper does not make.]
  \item[D-AIFILL] See \Cref{sec:ai-disclosure}. [FILL: name specifically
    which passages, tables, or code were AI-drafted vs. author-verified.
    Every field is stated explicitly as "none identified", never omitted.]
  \item[D-TIER] [FILL: name the single highest tier claimed anywhere in this
    paper (\Cref{tab:tiers}) and the one reason it cannot go higher yet.]
  \item[D-INDEPENDENCE] [FILL: name the independence class (\(I0\)--\(I5\):
    self / same-model AI / same-vendor other model / decorrelated-vendor AI /
    mechanical or original-record check / external human) of every check
    this paper's claims received, per claim-card (\Cref{sec:five-questions}).
    This paper does not seek, and does not treat as available,
    any external or institutional validation as a lever that would make its
    claims more true on its own say-so; an outside reader's input, if used,
    is optional-level evidence at the independence level it actually reaches,
    never a certification gate. Same-model self-approval is \(I0\) and is not
    counted as an independent check.]
  \item[D-DVP-NOT-K2] A decorrelated-vendor AI review pass (DVP) is at most
    an \(I3\) route. \textbf{DVP $\neq$ K2}: no combination of AI-only
    checks, however many vendors, raises a claim to the K2 (externally
    certified) knowledge state (\Cref{sec:kstate}) --- only a named, class-5
    independent human check does. [FILL: state this paper's actual K-state
    and the independence class each load-bearing claim reached, or: no
    claim in this paper is asserted above K1.]
  \item[D-LEGAL-NEQ-EPISTEMIC] [FILL, or: none identified.] Any license,
    permit, certification, or regulatory approval referenced in this paper
    makes an action lawful; it does not raise this paper's epistemic warrant
    and does not certify any claim as true.
  \item[D-NOT-DIAGNOSTIC] [FILL if health/clinical-adjacent, else: not
    applicable --- this paper makes no health, clinical, or treatment
    recommendation.]
  \item[D-TRANSLATION] [FILL if a Thai-language companion exists: missing /
    machine / reviewed / verified, per record. Else: not applicable ---
    single-language artifact.]
  \item[D-DISSENT-PRESERVED] [FILL: pointer to any dissent or deviation
    record connected to this paper's claims, or: none recorded to date. A
    dissent record is never silently deleted by a later draft.]
  \item[D-COMPARISON] This paper does not claim priority or precedence over
    other work. \Cref{sec:related} answers exactly three questions --- which
    problem this paper solves, by which method, and how the nearest
    neighbours do it the same or differently (\Cref{tab:neighbours}) --- and
    every neighbour relation is stated as \emph{same}, \emph{different}, or
    \emph{cited}, never as a priority contest. "Adopted from" is written only
    where a human instructed the adoption and a Blackbox Note line or
    decision id names it (\Cref{sec:blackbox}).
  \item[D-CITATION-UNVERIFIED] [FILL, or: none identified.] Every citation in
    \Cref{sec:related} and elsewhere carries a \texttt{fetch\_status}
    (\texttt{FETCHED\_VERIFIED} / \texttt{NOT\_FETCHED} / \texttt{FETCH\_
    FAILED} / \texttt{MISMATCH} / \texttt{CHALLENGED} / \texttt{SUPERSEDED}).
    Source existence does not, by itself, mean the cited source supports the
    claim attached to it (metadata verification $\neq$ claim-match
    verification). Any row without a verified citation card is tiered
    \OpenTier{}, not upgraded on the strength of the citation alone.
  \item[D-BLACKBOX-NOTE] This paper carries the mandatory Blackbox Note
    appendix (\Cref{sec:blackbox}, "how this work was made"): the curated
    verbatim human lines that produced this paper's claims and the cooking
    log of every transformation applied to them. [FILL: Blackbox Note id(s),
    e.g.\ \texttt{BB-YYYY-MM-DD-NN}.]
  \item[D-K-STATE] [FILL: state this paper's overall knowledge state
    (\Cref{sec:kstate}): \texttt{K0} internal working note, \texttt{K1}
    public-provisional (published, independence class named per claim, no
    class-5 human check yet), \texttt{K2} externally certified (a named
    class-5 human independently checked the load-bearing claims), or
    \texttt{K3} superseded/withdrawn. A paper may mix K-states across claims;
    if so, name the K-state per claim-card row in \Cref{tab:claim-matrix}
    instead of one blanket value.]
  \item[D-AUTHORSHIP] See \Cref{sec:ai-disclosure}. The research direction,
    standpoint, and final public commitment are the human author's; AI
    contribution is disclosed per claim (\texttt{produced\_by},
    \Cref{tab:claim-matrix}) and never substitutes for the author's own
    judgment or responsibility (\textbf{AIContribution $\neq$
    EpistemicResponsibility}).
\end{disclaimers}
